\documentclass[12pt]{article}
% AAU theme adapted from the preamble supplied by the project team.
\usepackage[a4paper,margin=1in]{geometry}
\usepackage{lmodern}
\usepackage[T1]{fontenc}
\usepackage{microtype}
\usepackage{graphicx}
\usepackage{float}
\usepackage{amsmath}
\usepackage[table]{xcolor}
\usepackage{titlesec}
\usepackage{fancyhdr}
\usepackage[most]{tcolorbox}
\PassOptionsToPackage{obeyspaces}{url}
\usepackage[hidelinks]{hyperref}
\usepackage{array}
\usepackage{tabularx}
\usepackage{enumitem}
\usepackage{booktabs}
% Vector UML artwork and multipage test tables.
\usepackage{longtable}
\usepackage{pdflscape}
\definecolor{aauRed}{RGB}{164,22,35}
\definecolor{aauGreen}{RGB}{0,122,88}
\definecolor{darkGray}{RGB}{60,60,60}
\definecolor{softGray}{RGB}{245,245,245}
\pagestyle{fancy}\fancyhf{}\fancyfoot[C]{\thepage}
\renewcommand{\headrulewidth}{0pt}\renewcommand{\footrulewidth}{0pt}
\titleformat{\section}{\normalfont\Large\bfseries\color{aauRed}}{\thesection}{0.6em}{}[{\color{aauRed!65}\titlerule[1.1pt]}]
\titlespacing*{\section}{0pt}{12pt}{8pt}
\titleformat{\subsection}{\normalfont\normalsize\bfseries\color{aauRed!85!black}}{\thesubsection}{0.6em}{}
\titlespacing*{\subsection}{0pt}{10pt}{4pt}
\newtcolorbox{objectivebox}{enhanced,colback=aauRed!5,colframe=aauRed!80!black,boxrule=0.3mm,arc=1mm,left=10pt,right=10pt,top=6pt,bottom=6pt,borderline west={2.5pt}{0pt}{aauRed}}
% Each short use/test case stays on one page, including its heading.
\newtcolorbox{contentbox}[1][]{enhanced,colback=softGray,colframe=black!35,boxrule=0.25mm,arc=2mm,left=9pt,right=9pt,top=7pt,bottom=7pt,before skip=9pt,after skip=9pt,title={\textbf{#1}},fonttitle=\bfseries,coltitle=darkGray,colbacktitle=black!8}
\newcommand{\blankline}{\rule{\linewidth}{0.35pt}}
\newcommand{\insertfigure}[3]{\begin{figure}[H]\centering\fbox{\parbox[c][#2][c]{0.88\textwidth}{\centering\textcolor{black!50}{\Large #1}}}\caption{#3}\end{figure}}
\newcommand{\screen}[3][0.98]{\begin{figure}[H]\centering\includegraphics[width=#1\linewidth]{../assets/#2}\caption{#3}\end{figure}}
% Give wide UML models a dedicated landscape page with readable labels.
\newenvironment{umlpage}{%
\newgeometry{margin=0.65in}\begin{landscape}\pagestyle{empty}%
\AddToHook{shipout/foreground}[arena/umlfooter]{%
\begin{tikzpicture}[remember picture,overlay]
\node[rotate=90,anchor=south,inner sep=0pt] at ([xshift=-23pt]current page.east)
{\normalfont\normalsize\thepage};
\end{tikzpicture}}%
}{\end{landscape}\RemoveFromHook{shipout/foreground}[arena/umlfooter]\pagestyle{fancy}\restoregeometry}
\newcommand{\umlfigure}[2]{\begin{figure}[H]\centering
\includegraphics[width=\linewidth,height=13cm,keepaspectratio]{../assets/uml/#1.pdf}
\caption{#2}\end{figure}}
\providecommand{\reportdate}{10 November 2026}
% Two intentionally sized figures fit together with a heading on an A4 page.
\newcommand{\screenpair}[5][0.30]{%
\begin{figure}[H]\centering
\includegraphics[width=0.96\linewidth,height=#1\textheight,keepaspectratio]{../assets/#2}\caption{#3}
\vspace{8pt}
\includegraphics[width=0.96\linewidth,height=#1\textheight,keepaspectratio]{../assets/#4}\caption{#5}
\end{figure}}
\setlength{\parindent}{0pt}\setlength{\parskip}{5pt}
\setlist[itemize]{leftmargin=1.5em,itemsep=2pt,topsep=3pt}
\setlist[enumerate]{leftmargin=1.7em,itemsep=3pt,topsep=4pt}
\setlength{\emergencystretch}{2em}
\widowpenalty=10000\clubpenalty=10000\displaywidowpenalty=10000
\brokenpenalty=10000
\raggedbottom
\renewcommand{\arraystretch}{1.16}
\setlength{\abovecaptionskip}{7pt}
\setlength{\belowcaptionskip}{4pt}
\setlength{\intextsep}{12pt plus 2pt minus 2pt}
\newcolumntype{Y}{>{\raggedright\arraybackslash}X}
\BeforeBeginEnvironment{tabularx}{\par\noindent}
\AfterEndEnvironment{tabularx}{\par}
\newcommand{\code}[1]{{\small\expandafter\nolinkurl\expandafter{\detokenize{#1}}}}
\newcommand{\command}[1]{\mbox{\small\ttfamily\detokenize{#1}}}
\newcommand{\projecttitle}[2]{%
\hypersetup{pageanchor=false}\begin{titlepage}\centering\vspace*{-0.45cm}
\noindent\begin{minipage}[c]{0.66\textwidth}\raggedright
{\Large\bfseries Al Ain University (AAU)\par}
\end{minipage}\hfill
\begin{minipage}[c]{0.28\textwidth}\raggedleft
% The supplied PNG is unchanged; trim removes only its transparent margins.
\includegraphics[width=3.1cm,trim=1731bp 871bp 1735bp 878bp,clip]{../../aau.png}
\end{minipage}\par
\vspace{0.45cm}
{\scshape\large Object Oriented Analysis and Design\par}
\vspace{0.10cm}
{\normalsize System Analysis \& Design\par}
\vspace{0.10cm}
{\small Course codes: 0103430 \;|\; 0102331\par}
\vspace{0.3cm}\noindent\rule{\linewidth}{0.5mm}\par\vspace{0.3cm}
{\huge\bfseries Arena\par}
\vspace{0.18cm}{\Large\color{aauRed}#1\par}
\vspace{0.15cm}{\normalsize\color{darkGray}#2\par}
\vspace{0.12cm}{\color{aauRed}\rule{\linewidth}{1.2pt}\par}
\vspace{0.35cm}{\normalsize\bfseries Prepared by\par}\vspace{0.15cm}
\begin{tabular}{@{}lr@{}}\textbf{Name}&\textbf{Student ID}\\\midrule
Khaled Alkhawaldeh&202410091\\Abdulrahman Alnaqbi&202340042\\Yousif Alhammadi&202310707\\Ahmed Alsamma&202320192\end{tabular}\par
\vspace{0.22cm}{\color{aauRed!55}\rule{0.6\textwidth}{1.2pt}\par}
\vspace{0.15cm}{\normalsize\textbf{Section:} 61D\par}
\vspace{0.15cm}{\normalsize\textbf{Submitted to:} Dr. Yazeed Ghadi\par}
\vspace{0.3cm}
\begin{tcolorbox}[colback=gray!5,colframe=black!70,title=\textbf{Project Summary},fonttitle=\bfseries,width=0.95\textwidth,boxrule=0.4mm,arc=3mm,left=7pt,right=7pt,top=4pt,bottom=4pt]
A Java desktop workspace for organizing campus competitions, coordinating teams, collecting submissions, evaluating work, and publishing results. Project, programming, and quiz scoring are separated from score-ranking, knockout, and league progression through object-oriented strategies. The application is developed for Al Ain University in the United Arab Emirates.
\end{tcolorbox}
\vfill{\normalsize Fall, 2026--2027\par}{\small Version 2.0 \;|\; \reportdate\par}
\end{titlepage}\pagenumbering{roman}\hypersetup{pageanchor=true}\tableofcontents\clearpage\pagenumbering{arabic}}

\renewcommand{\reportdate}{10 November 2026}
\begin{document}
\projecttitle{System Analysis}{Requirements, domain models, and interaction design}
\section{Introduction}
\subsection{Project Background and Problem Statement}
Campus competitions bring together students, organizers, and judges, but a collection of registration forms and spreadsheets does not express how different competitions actually operate. A hackathon needs a project submission and independent rubric reviews. A programming contest needs an official solved-problem count and a time penalty. A robotics tournament needs opponents, winners, and progression through rounds.

\textbf{Arena} provides one coherent workspace for those variations. The instructor accepted the idea and specifically requested a scope beyond basic registration and scheduling. This analysis therefore treats format-dependent rules, guarded workflows, and fair result calculation as the central problem.
\subsection{Project Objectives}
\begin{objectivebox}
Deliver a runnable Java application that coordinates teams, validates entries, collects work, separates competition types from tournament formats, and publishes reproducible results. Demonstrate encapsulation, immutable domain objects, polymorphic rule strategies, explicit lifecycle states, and separation of interface, application logic, and persistence.
\end{objectivebox}
The implementation must reject invalid actions in the application service even when the graphical interface is bypassed. Its reports must describe the delivered implementation and distinguish executed tests from future validation.
\subsection{Project Scope}
Included: local account management; four roles; team creation and self-service membership; captain-led registration and withdrawal; configurable draft rules; submission links and descriptions; assigned judging; project, programming, quiz, and general-match types; score ranking, knockout, and round-robin formats; publication; CSV export; persistent storage; audit history; and supporting UML and interface artifacts.

Excluded from this release: institutional single sign-on, email delivery, cloud synchronization, automatic code execution, payment processing, venue reservations, file uploads, automatic disqualification, and a production deployment. Competition deadlines and bracket rounds are supported; a calendar with resource-conflict scheduling is not implemented.
\subsection{Proposed Technologies}
Java 17 language/API compatibility, Java Swing/AWT for the desktop interface, Java NIO for local storage and file locking, Java cryptography for salted password derivation, and LaTeX with vector UML illustrations for documentation. A portable JDK 21 was used to compile with \command{--release 17}. Build and test scripts require no external application libraries.
\subsection{Team Members}
\begin{tabularx}{\linewidth}{Yr}\toprule Name & Student ID\\\midrule
Khaled Alkhawaldeh & 202410091\\ Abdulrahman Alnaqbi & 202340042\\ Yousif Alhammadi & 202310707\\ Ahmed Alsamma & 202320192\\\bottomrule\end{tabularx}
Course names: Object Oriented Analysis and Design / System Analysis \& Design. Course codes: 0103430 / 0102331. Section: 61D. Instructor: Dr. Yazeed Ghadi.

\section{Overview}
\subsection{System Description}
A user signs into a single-machine competition workspace. Organizers configure draft competitions, assign judges, open registration, and advance the lifecycle. Students form teams and captains enter competitions. Depending on the selected format, work is submitted and evaluated, or matches are decided and winners advance. Published results are visible to every authenticated role.
\subsection{Target Users and Stakeholders}
\begin{tabularx}{\linewidth}{lY}\toprule Stakeholder & Interest and responsibility\\\midrule
Student & Find events, collaborate, enter as a team, submit work, and see official results.\\
Captain & A student who created a team; registers, withdraws, and submits for that team.\\
Organizer & Own competitions, set rules and judges, enforce readiness, and publish results.\\
Judge & Review only assigned competitions and record scores or match winners.\\
Administrator & Create accounts and manage any competition; inspect audit history.\\
Instructor & Assess completeness and the connection between models, code, and tests.\\\bottomrule\end{tabularx}
\subsection{System Features}
\begin{tabularx}{\linewidth}{lY}\toprule Capability & Delivered behavior\\\midrule
Project ranking & Submission; 40/35/25 rubric; mean of all assigned reviews.\\
Programming ranking & Submission; solved count first, then lower penalty; one official result.\\
Quiz ranking & Points out of a configured maximum; no submission; equal points share rank.\\
Knockout & General match or quiz; 2--16 entrants with byes; automatic next rounds.\\
Round robin & General match or quiz; every pair once; 3/1/0 points and deterministic ties.\\\bottomrule\end{tabularx}
\subsection{Assumptions and Constraints}
The application is used by one process on one trusted computer. Different people sign in sequentially; it is not a multi-client server. Local operating-system access controls protect the data file. Competition deadlines use the computer's local time; documents use the team's UAE context. Roster membership is voluntary self-joining and is frozen once a team has any registration. Team size is a maximum, not an exact minimum. Each competition needs at least two teams to start.

\section{Functional Requirements}
\subsection{User Management}
\textbf{FR-01 Authentication.} The system shall verify an email/password pair, support sign-out, and clear an existing session after a failed login. Email matching shall be case insensitive.\par
\textbf{FR-02 Accounts and credentials.} Only administrators shall create accounts and assign roles. Email addresses shall be unique, passwords shall have at least eight characters, and authenticated users shall be able to change their own password after verifying the old password.
\subsection{Competition Management}
\textbf{FR-03 Competition lifecycle.} Organizers and administrators shall create draft competitions with title, description, competition type, compatible tournament format, quiz maximum where relevant, maximum entries, maximum team size, deadline, and assigned judges. Only the owning organizer or an administrator shall edit or advance a competition. Rule editing shall stop when registration opens. Stages shall advance in order: Draft, Registration, Active, Judging, Published, Archived.

\textbf{Acceptance guards:} registration must open before the deadline; starting requires at least two entries; project/programming ranking events must start before their submission deadline; entering Judging requires all required submissions or all match fixtures completed. Quiz ranking requires no submission. Publication requires every assigned project review or one official programming/quiz result per entry; match formats must already be certified.
\subsection{Registration and Team Management}
\textbf{FR-04 Entry validation.} Only a team captain shall register or withdraw a team during registration. The system shall reject a full competition, an oversized team, repeated team registration, or a student appearing in two teams within the same competition. Entry after the deadline shall be rejected.\par
\textbf{FR-05 Team membership.} Students shall create uniquely named teams, join unregistered teams, and leave teams they do not captain. A team shall contain at most eight members. Registered rosters shall be locked to preserve eligibility. Captains cannot leave their own team in this release.
\subsection{Submission Management}
\textbf{FR-06 Submissions.} Captains of registered teams shall submit a title, description, and HTTP(S) project link during the Active phase and strictly before the deadline. Resubmission shall replace that entry's earlier submission. Quiz and all match formats shall reject submissions. Students shall only see submissions belonging to their own teams; assigned judges and competition managers shall see all entries in their competition.
\subsection{Judging and Evaluation}
\textbf{FR-07 Assigned review.} Only assigned judges shall record evaluations during Judging. The system shall require a submission for project/programming ranking only, score ranges appropriate to the type, and nonempty feedback. A judge may revise their own score before publication. Published scores shall be immutable.\par
\textbf{FR-09 Programming results.} A result shall contain solved problems from 0--10, elapsed minutes from 0--300, and wrong attempts from 0--100. One assigned judge owns the official result for each team; a second judge cannot replace it. No submitted code is executed.\par
\textbf{FR-10 Knockout workflow.} Starting shall require 2--16 registered teams, seeded by registration order. The opening bracket is rounded up to the next power of two; early seeds receive automatic first-round byes. Assigned judges shall record winners only while Active. A winner must be a participant in that match. General knockout events use winner selection; quiz knockout events use bounded points for both teams and reject draws until an external tie-break produces a decisive result. Results are final; when every match in a round is complete, winners are paired for the next round.
\subsection{Results and Leaderboard}
\textbf{FR-08 Results and export.} The system shall publish rankings only after the relevant completeness checks. Organizers and administrators may preview provisional standings; other roles wait for publication. Equal score-ranking results share a rank, using competition ranking (1, 1, 3). Knockout ranks use the round reached, with the champion first. League ranks follow their deterministic tie-break order. Export shall produce CSV and neutralize leading spreadsheet formula characters in text cells.\par
\textbf{FR-11 Audit.} Successful changes shall record actor, local timestamp, and action. Administrators and organizers shall see the most recent 2,000 events. Failed operations shall not appear as successful events.

\subsection*{Scoring Policies}
For a project scored by judge $j$, with three scores between 0 and 100:
\[S_j=0.40I_j+0.35E_j+0.25P_j,\qquad S_{team}=\frac{1}{n}\sum_{j=1}^{n}S_j.\]
All assigned judges must score all entries before publication. Scores are averaged without rounding internally and displayed to two decimal places. Equality uses a $10^{-6}$ tolerance.

For programming, sort solved count descending and penalty ascending:
\[\mathrm{penalty}=\mathrm{elapsed\ minutes}+20\times\mathrm{wrong\ attempts}.\]
This is an intentionally simplified coursework policy, not a claim to implement all ICPC rules. Tied solved counts and penalties share a rank. Knockout standings rank the champion first and then use the round reached; teams eliminated at the same depth share rank. The displayed score counts played-match wins and excludes byes, so a bye cannot distort the champion's rank.

\subsection{Arena Extensions: Competition Type and Tournament Format}
The following extension requirements record the requested Arena update. They supplement the original academic deliverables; they do not replace the assignment instructions.

\textbf{FR-12 Quiz scoring.} A quiz ranking event shall record one official integer points result per registered team, from zero to a configured maximum of 1--10,000. Extra scoring inputs shall be zero. The owning judge may revise that result while Judging. Equal points share a rank and tied first-place entries share the win. Questions and any knockout tie-break are administered outside the application.

\textbf{FR-13 Round-robin league.} Starting shall generate exactly $n(n-1)/2$ fixtures for $n=2,\ldots,16$ teams. Each pair meets once and no team appears twice in one round; odd counts include rest rounds. Assigned judges record nonnegative integer scores, bounded by the quiz maximum or 10,000 for general matches. Wins award 3 points, draws 1, losses 0. Rank by points, wins, score difference, and scores for (descending), then case-insensitive team name and stable team ID (ascending). Every fixture must finish before certification and publication. Confirmed match results cannot be revised.

\textbf{FR-14 Compatible configurations.} Creation and draft editing shall enforce this matrix in the service, and the UI shall offer only supported choices:
\begin{tabularx}{\linewidth}{Yccc}\toprule Type & Score ranking & Knockout & Round robin\\\midrule
Project & Yes & No & No\\ Programming & Yes & No & No\\ Quiz / academic & Yes & Yes & Yes\\ General match & No & Yes & Yes\\\bottomrule\end{tabularx}
Project and programming retain their submission and review workflows. Quiz requires no project submission. A general-match type preserves the original generic knockout functionality.

\section{Nonfunctional Requirements}
\subsection{Performance}
\textbf{NFR-01:} interactive operation is intended for small campus datasets (up to 64 teams per score-ranking competition and 16 per match tournament). A target of under one second for routine navigation is a design goal, not a measured benchmark. Password derivation and disk writes can briefly block the Swing thread; large-scale performance testing remains future work.
\subsection{Security}
\textbf{NFR-02:} passwords shall be stored as salted PBKDF2-HMAC-SHA256 derivations, never plaintext; service methods shall enforce roles and ownership. Serialized reads shall use class and size limits. The local file is not encrypted and protection against a malicious user who can modify it is outside scope.
\subsection{Usability}
\textbf{NFR-03:} visible phase indicators, consistent navigation, labeled controls, clear empty states, and actionable validation errors shall support users without technical knowledge. The application shall support a minimum 1100 by 720 window. Visual checks at the default size are recorded; formal accessibility and participant usability studies have not been performed.
\subsection{Reliability}
\textbf{NFR-04:} successful updates shall persist; failed saves shall restore the in-memory state; an exclusive file lock shall reject a second process. Corrupt data shall trigger an error rather than silently reset. Atomic replacement is used where the filesystem supports it, with normal replacement as fallback.
\subsection{Maintainability}
\textbf{NFR-05:} business policies shall be testable without the UI. Strategies shall encapsulate format validation and scoring. Requirements shall map to executable checks.

\begin{umlpage}
\section{System Models}
\subsection{Use Case Model}
\umlfigure{01_use_case}{Actor and use-case overview. All four roles share signed-in user capabilities. Authentication is a precondition for protected use cases.}
\end{umlpage}
\begin{contentbox}[UC-01: Register a team]
\textbf{Actor:} student captain. \textbf{Preconditions:} signed in, team exists, registration open, deadline not reached.\\
\textbf{Main flow:} choose competition; choose own team; service verifies captain, phase, capacity, size, and overlapping members; store entry and audit event; refresh teams.\\
\textbf{Alternatives:} duplicate, full, oversized, overlapping, or late entry produces an error and no mutation. A save failure restores the previous state.\\
\textbf{Postcondition:} one registration exists and the roster is locked.
\end{contentbox}
\begin{contentbox}[UC-02: Evaluate a project]
\textbf{Actor:} assigned judge. \textbf{Preconditions:} Judging phase; registered entry has submitted.\\
\textbf{Main flow:} read submission; select team; enter three criteria and feedback; validate ranges; replace only this judge's previous evaluation; save.\\
\textbf{Alternatives:} unassigned judge, missing work, invalid score, or wrong phase is rejected.\\
\textbf{Postcondition:} exactly one evaluation for this judge and entry; results remain private until publication.
\end{contentbox}
\begin{contentbox}[UC-03: Decide a knockout match]
\textbf{Actor:} assigned judge. \textbf{Preconditions:} Active knockout; undecided match.\\
\textbf{Main flow:} choose one competing team as winner; confirm the final result; save winner; if the round is complete, pair winners into the next round.\\
\textbf{Alternatives:} unrelated team, unassigned judge, already decided match, or wrong phase is rejected.\\
\textbf{Postcondition:} winner is immutable; at most one next round is generated.
\end{contentbox}
\begin{contentbox}[UC-04: Publish results]
\textbf{Actor:} owner or administrator. \textbf{Preconditions:} Judging phase.\\
\textbf{Main flow:} request publication; verify every required review or certified bracket; advance phase; display official standings to all signed-in users.\\
\textbf{Alternatives:} incomplete reviews leave the competition in Judging.\\
\textbf{Postcondition:} published evaluations cannot be modified; CSV export is available.
\end{contentbox}
\begin{umlpage}
\subsection{Dynamic Models}
\umlfigure{02_registration_sequence}{Successful team registration. Validation failure returns before mutation; persistence failure restores the snapshot.}
\clearpage\subsubsection*{Competition Lifecycle}
\umlfigure{03_competition_state}{Competition state machine and transition guards.}
\clearpage\subsubsection*{Format-Specific Workflow}
\umlfigure{04_competition_activity}{Activity overview separating score-ranking, knockout, and league workflows.}
\end{umlpage}
\begin{umlpage}
\subsection{Object and Class Model}
\umlfigure{05_conceptual_domain}{Analysis class relationships. Stored relationships use stable identifiers. Each captain belongs to the team's members.}
\end{umlpage}
\clearpage\subsection{User Interface}
The primary navigation is Sign in $\rightarrow$ Overview $\rightarrow$ Competitions $\rightarrow$ Competition workspace. A workspace contains Teams, relevant Submissions/Reviews, Bracket or Fixtures, and Leaderboard. My teams and Results are global destinations; People is administrator-only and Activity is available to organizers and administrators. Role-specific action buttons reflect, but do not replace, service authorization.
The separate Interface Mock-ups document includes the complete wireframe set, seven browser-prototype captures and thirteen rendered application screens and dialogs.
\screenpair{prototype/01-overview.png}{Clickable browser prototype dashboard. Demo data resets on refresh.}{screenshots/04-competition.png}{Implemented competition workspace with format policy and lifecycle indicator.}

\clearpage\section{Glossary \& References}
\begin{tabularx}{\linewidth}{lY}\toprule Term & Meaning\\\midrule
Entry & One team's registration in one competition.\\
Captain & Student responsible for registering and submitting for a team.\\
Strategy & Interchangeable object encapsulating format-specific behavior.\\
Phase guard & Condition checked before the lifecycle advances.\\
Official result & The single judge-owned programming or quiz ranking score for a team.\\
Provisional standings & Manager-only rankings before publication.\\
Atomic replacement & Replacing the saved file in one filesystem operation where supported.\\\bottomrule\end{tabularx}
\begingroup\raggedright\begin{enumerate}
\item Object Management Group. \emph{Unified Modeling Language, version 2.5.1}. \url{https://www.omg.org/spec/UML/2.5.1}. Accessed 29 October 2026. Used as the notation reference for the UML models.
\item Oracle. \emph{Java SE 17: javax.swing package}. \url{https://docs.oracle.com/en/java/javase/17/docs/api/java.desktop/javax/swing/package-summary.html}. Accessed 29 October 2026. Used for desktop components and event-dispatch-thread conventions.
\item Oracle. \emph{Java SE 17: ObjectInputFilter}. \url{https://docs.oracle.com/en/java/javase/17/docs/api/java.base/java/io/ObjectInputFilter.html}. Accessed 29 October 2026. Used for bounded, filtered repository deserialization.
\item Project materials. Task outlines in \code{requirements/system-analysis-requirements.txt}, \code{requirements/system-design-requirements.txt}, \code{requirements/implementation-report-requirements.txt}, \code{requirements/testing-requirements.txt}, and \code{requirements/sample-data-requirements.txt}; instructor feedback supplied by the team. These define the coursework scope, not an independent evaluation of the implementation.
\end{enumerate}\endgroup

\end{document}
